Kurs:Algebraische Kurven (Osnabrück 2012)/Arbeitsblatt 28/latex
Kurs:Algebraische Kurven (Osnabrück 2012)/Arbeitsblatt 28/latex


\setcounter{section}{28}


  \zwischenueberschrift{Aufwärmaufgaben}


\inputaufgabe

{}

{


Es sei


\mavergleichskette

{\vergleichskette

{P
}

{ = }{ \left( a_0   , \,   \ldots  , \,   a_n                 \right)
}

{ \in }{ {\mathbb P}^{ n }_{ K}
}

{    }{
}

{   }{
}

}
{}{}{}
ein Punkt im
\definitionsverweis {projektiven Raum}{}{.}
Zeige, dass es eine offene affine Umgebung


\mavergleichskette

{\vergleichskette

{U
}

{ \cong }{ { {\mathbb A}_{  K }^{  n   } }
}

{ \subset }{ {\mathbb P}^{ n }_{ K}
}

{    }{
}

{   }{
}

}
{}{}{}
derart gibt, dass $P$ in diesem affinen Raum dem Nullpunkt entspricht.


}

{} {}


\inputaufgabe

{}

{


Es sei


\mavergleichskette

{\vergleichskette

{ D_+(L)
}

{ \cong }{ { {\mathbb A}_{  K }^{  n   } }
}

{ \subseteq }{ {\mathbb P}^{ n }_{ K}
}

{    }{
}

{   }{
}

}
{}{}{,}
wobei $L$ eine homogene Linearform im zugehörigen Polynomring

\mathl{K[X_0 , \ldots , X_n]}{} sei. Zeige, dass die Zariski-Topologie auf dem projektiven Raum die Zariski-Topologie auf dem affinen Raum induziert.


}

{} {}


\inputaufgabe

{}

{


Definiere zu jedem


\mavergleichskette

{\vergleichskette

{ n
}

{ \in }{ \Z
}

{  }{
}

{    }{
}

{   }{
}

}
{}{}{}
das Potenzieren

\mathl{x \mapsto x^n}{} als Morphismus der projektiven Gerade auf sich selbst. Wie sehen die Fasern unter diesem Morphismus aus?


}

{} {}


\inputaufgabe

{}

{


Bestimme den
\definitionsverweis {projektiven Abschluss}{}{} der durch


\mathdisp {V \left(  \left(  X^2+Y^2  \right)^2- 2X \left(  X^2+Y^2  \right) -Y^2  \right)} {  }


gegebenen

\definitionswortenp{Kardioide}{} über den komplexen Zahlen und insbesondere die \anfuehrung{Punkte im Unendlichen}{.}


}

{} {}


\inputaufgabe

{}

{


Zeige, dass die
\definitionsverweis {Kegelabbildung}{}{}
\maabbdisp {} { { {\mathbb A}_{  K  }^{ n+1  } } \setminus \{0\} } { {\mathbb P}^{ n }_{ K }
} {}
ein Morphismus von quasiprojektiven Varietäten ist.


}

{} {}


\inputaufgabe

{}

{


Zeige durch ein Beispiel, dass die
\definitionsverweis {Kegelabbildung}{}{}
\maabbdisp {} { { {\mathbb A}_{  K  }^{ n+1  } } \setminus \{0\} } { {\mathbb P}^{ n }_{ K }
} {}
nicht
\definitionsverweis {abgeschlossen}{}{}
sein muss.


}

{} {}


\inputaufgabe

{}

{


Es seien $X$ und $Y$
\definitionsverweis {quasiprojektive Varietäten}{}{}
und sei
\maabb {\varphi} { X } { Y
} {}
eine
\definitionsverweis {stetige Abbildung}{}{.}
Es sei


\mavergleichskette

{\vergleichskette

{ Y
}

{ = }{ \bigcup_{i \in I} U_i
}

{  }{
}

{    }{
}

{   }{
}

}
{}{}{}
eine
\definitionsverweis {offene Überdeckung}{}{.}
Zeige, dass $\varphi$ genau dann ein
\definitionsverweis {Morphismus}{}{}
ist, wenn die Einschrän\-kungen
\maabbdisp {\varphi_i} { \varphi^{-1} (U_i)  } { U_i
} {}
für jedes $i$ Morphismen sind


}

{} {}


\inputaufgabe

{}

{


Es sei $K$ ein Körper. Bestimme den globalen Schnittring


\mathdisp {\Gamma   \left(   {\mathbb P}^{1}_{K} ,  {\mathcal O}_{  {\mathbb P}^{1}_{K}  }   \right)} { . }


Was folgt daraus für einen Morphismus
\maabb {} { {\mathbb P}^{1}_{K} } { {\mathbb A}^{1}_{K}
} {?}


}

{} {}


\inputaufgabe

{}

{


Man definiere und charakterisiere, wann eine irreduzible quasiprojektive Varietät \stichwort {normal} {} ist.


}

{} {}


\inputaufgabegibtloesung

{}

{


Es sei $K$ ein
\definitionsverweis {Körper}{}{.}
Betrachte die affine ebene Kurve


\mavergleichskettedisp

{\vergleichskette

{C
}

{ =} { V \left(  Y-X^3+X+2  \right)
}

{ } {
}

{ } {
}

{ } {
}

}
{}{}{.}
Definiere einen Isomorphismus zwischen $C$ und der affinen Geraden ${\mathbb A}^{1}_{ K }$. Lässt sich ein solcher Isomorphismus zu einem Isomorphismus zwischen ${\mathbb P}^{1}_{K}$ und dem
\definitionsverweis {projektiven Abschluss}{}{}


\mavergleichskette

{\vergleichskette

{ \overline{C}
}

{ \subset }{ {\mathbb P}^{2}_{K}
}

{  }{
}

{    }{
}

{   }{
}

}
{}{}{}
fortsetzen?


}

{} {}


  \zwischenueberschrift{Aufgaben zum Abgeben}


\inputaufgabe

{}

{


Es sei


\mavergleichskette

{\vergleichskette

{ {\mathbb K}
}

{ = }{ \R
}

{  }{
}

{    }{
}

{   }{
}

}
{}{}{}
oder $= {\mathbb C}$. Es sei


\mavergleichskette

{\vergleichskette

{ H
}

{ \subset }{ {\mathbb K}^{n+1}
}

{  }{
}

{    }{
}

{   }{
}

}
{}{}{}
ein $n$-dimensionaler
\definitionsverweis {affiner Unterraum}{}{,}
der den Nullpunkt nicht enthält, und es sei $\tilde{H}$ der dazu parallele Unterraum durch den Nullpunkt. Es sei


\mavergleichskette

{\vergleichskette

{ U
}

{ \subseteq }{ H
}

{  }{
}

{    }{
}

{   }{
}

}
{}{}{}
eine in


\mavergleichskette

{\vergleichskette

{ H
}

{ \cong }{ {\mathbb K}^n
}

{  }{
}

{    }{
}

{   }{
}

}
{}{}{}
offene Menge
\zusatzklammer {in der metrischen Topologie} {} {}
und es sei $V$ die Vereinigung aller Geraden durch den Nullpunkt und durch einen Punkt von $U$. Zeige, dass der Durchschnitt von $V$ mit

\mathl{{\mathbb K}^{n+1} \setminus \tilde{H}}{} offen ist.


}

{} {}


\inputaufgabe

{}

{


Bestimme für die
\definitionsverweis {Kegelabbildung}{}{}
\maabbdisp {} { { {\mathbb A}_{  K  }^{ n+1  } } \setminus \{0\}} { {\mathbb P}^{ n }_{ K }
} {}
den Zariski-Abschluss im ${\mathbb P}^{ n }_{ K }$ des Bildes einer abgeschlossenen Menge

\mathl{V( {\mathfrak a}) \cap { {\mathbb A}_{  K  }^{ n+1  } } \setminus \{0\}}{.}


}

{} {}


\inputaufgabe

{}

{


Es sei $X$ eine irreduzible
\definitionsverweis {quasiprojektive Varietät}{}{}
mit Funktionenkörper


\mavergleichskette

{\vergleichskette

{ L
}

{ = }{ K(X)
}

{  }{
}

{    }{
}

{   }{
}

}
{}{}{.}
Es seien $U$ und

\mathbed {U_i} {}

 {i \in I} {}

 {} {} {} {,} offene Teilmengen mit


\mavergleichskette

{\vergleichskette

{ U
}

{ = }{ \bigcup_{i \in I} U_i
}

{  }{
}

{    }{
}

{   }{
}

}
{}{}{.}
Zeige, dass


\mavergleichskettedisp

{\vergleichskette

{ \Gamma (U, {\mathcal O} )
}

{ =} { \bigcap_{i \in I} \Gamma (U_i , {\mathcal O} )
}

{ } {
}

{ } {
}

{ } {
}

}
{}{}{}
ist, wobei der Durchschnitt in $L$ genommen wird.


}

{} {}


\inputaufgabe

{}

{


Es sei $K$ ein
\definitionsverweis {Körper}{}{}
und ${\mathbb P}^{ n }_{ K}$ der
\definitionsverweis {projektive Raum}{}{}
über $K$. Zeige, dass die Konstanten die einzigen globalen
\definitionsverweis {algebraischen Funktionen}{}{}
sind, d.h. es gilt


\mavergleichskettedisp

{\vergleichskette

{ \Gamma   \left(   {\mathbb P}^{ n }_{ K} ,  {\mathcal O}_{  {\mathbb P}^{ n }_{ K}  }   \right)
}

{ =} { K
}

{ } {
}

{ } {
}

{ } {
}

}
{}{}{.}


}

{Bemerkung: Diese Aussage gilt für jede zusammenhängende projektive Varietät über einem algebraisch abgeschlossenen Körper.} {}


<< | Kurs:Algebraische Kurven (Osnabrück 2012) | >>


PDF-Version dieses Arbeitsblattes


 Zur Vorlesung
(PDF)
